\begin{center}
	\LARGE\bfseries{РЕФЕРАТ}
\end{center}

Расчетно-пояснительная записка \pageref{LastPage} с., \totalfigures{} рис., \totaltables{} таблица, 8 источников, 2 приложения.

БАЗА ДАННЫХ, POSTGRESQL, ИГРОВЫЕ МОДИФИКАЦИИ, РЕЛЯЦИОННАЯ МОДЕЛЬ ДАННЫХ, ИНДЕКСИРОВАНИЕ, РОЛЕВАЯ МОДЕЛЬ, REDIS.

Расчетно-пояснительная записка посвящена разработке базы данных для платформы игровых модификаций. Целью работы является создание структуры хранения данных, обеспечивающей управление, поиск и хранение модификаций, а также поддержку разграничения прав доступа пользователей.

В ходе выполнения работы был проведен анализ предметной области, выделены основные сущности и их взаимосвязи, рассмотрены существующие аналоги и сформулированы требования к разрабатываемой системе. Выполнено сравнение различных моделей хранения данных и обоснован выбор реляционной модели в качестве основного решения.

На этапе проектирования разработана структура базы данных, построена диаграмма сущность--связь и схема отношений, выполнена нормализация данных до третьей нормальной формы. Также спроектирована ролевая модель доступа и разработан алгоритм работы триггера, обеспечивающего соблюдение бизнес-логики системы.

В технологической части база данных реализована средствами PostgreSQL. Созданы таблицы, ограничения целостности, индексы и триггеры. Для повышения производительности системы предусмотрено использование Redis в качестве кэширующего хранилища. Проведено тестирование разработанных механизмов.

В исследовательской части выполнен анализ влияния различных типов индексов PostgreSQL на производительность поиска. Исследованы индексы B-Tree, Hash, GIN, GiST и BRIN. По результатам исследования установлено, что индекс B-Tree является предпочтительным решением для поиска модификаций, поскольку использует наименьшее количество памяти и сопоставим по скорости с Hash индексом. Индекс BRIN обеспечивает наилучшие результаты среди других индексов при выполнении диапазонных запросов.